\chapter{Modelos matemáticos de sistemas nolineales}

\section{Generalidades}

Nolinealidad geométrica o del material.

En el segundo caso, histeresis ocurre cuando la fuerza es función de la historia de deformación (dependiente de la memoria del componente o material).

\section{Ecuaciones de movimiento nolineales}


\section{Modelos fenomenológicos}

\textcite{wang2023modeldrivendatadriven}

\subsection{Características}

Modelos con curvas esqueleto (backbone): polinomial y polinomial por trozos (piecewise)

Modelos basados en ecuaciones diferenciales

\subsection{Modelo elastoplástico}

Ref: Simo \& Hughes

Consideremos un modelo uniaxial elastico -- perfectamente plástico.

$$\sigma = E (\varepsilon - \varepsilon_p)$$

$$f = |\sigma| - \sigma_y \leq 0$$

$$\dot{\varepsilon}_p = \gamma \text{sgn}(\sigma)$$

Si $f < 0$ entonces $\gamma = 0$

Si $f = 0$ entonces $\gamma > 0$

Algoritmo \emph{return mapping}, donde $\varepsilon_p$ es una variable interna con la historia del proceso de plastificación.

\subsection{Modelo Bouc-Wen}

Ref: \textcite{song2006generalizedbouc}

Consideremos un elemento histerético con rigidez elástica $k_0$, desplazamiento de fluencia $d_y$, y razón de rigidez post-fluencia vs elástica $\alpha$. El modelo Bouc-Wen establece una relación nolineal controlado por la variable interna $z(t)$:

\begin{equation}
    f = \alpha k_0 d + (1-\alpha) k_0 d_y z
\end{equation}

% \begin{equation}
%     \dot{z} = \frac{1}{d_y} \left( A \dot{d} - \beta \dot{d} \left|z\right|^{n} - \gamma \left|\dot{d}\right| \left|z\right|^{n-1} \right)
% \end{equation}

\begin{equation}
    \dot{z} = \frac{1}{d_y} \dot{d} \left[ A  - \left|z\right|^{n} \left( \gamma + \beta \text{sgn}(\dot{d} z) \right)\right]
\end{equation}

\noindent donde $A$ y $n$ son parámetros que controlan la escala y la agudeza (\emph{sharpness}) de la curva de histeresis; $\beta$ y $\gamma$ controlan la forma de la curva de histeresis.

% En la literatura se recomienda $A=1$ y $\beta + \gamma = 1$.

\begin{lstlisting}[style=matlab-editor]
alpha = 0.;
beta = 0.55;
gamma = 0.45;
n = 2;
A = 1.;
\end{lstlisting}

\subsection{Modelo Duhem}

Utilizado para caracterizar sistemas friccionales.

\section{Modelos Hammerstein-Wiener}



%%%
% \section{Implementación computacional}
%%%

\section{Linealización equivalente}

\textcite{roberts2003randomvibration}